\paragraph{Minimal ethical risk.}
This work poses minimal ethical risk.
All experiments use synthetic tasks designed by the authors; no human subjects were involved and no personally identifiable information was collected or processed.
The LLM's own assessments serve as the reward signal for bandit learning; no human annotators or crowd workers were employed.
The system generates only task-completion responses (e.g., summaries, analyses) with no capacity for harmful content production.

\paragraph{Compute and environmental impact.}
Total compute for both experiments was approximately 52 hours on a Mac Mini M4 (Apple Silicon, ${\sim}$20W TDP).
LLM inference used the MiniMax M2.5 commercial API; local embedding inference used Qwen3-Embedding-0.6B.
The energy footprint is modest: under 2~kWh for local computation across both experiments, negligible relative to model training runs.
API-side energy consumption is not disclosed by the provider but is bounded by the small scale of inference (${\sim}$2{,}200 episodes, ${\sim}$4{,}000 tokens per episode).

\paragraph{Broader impact.}
The primary societal implication of this work is efficiency: reducing token consumption in LLM systems lowers inference costs and energy use at deployment time.
The input-assessment reward design, by avoiding output-level self-assessment, may contribute to safer feedback loops in deployed agentic systems by reducing the risk of reward hacking.
However, the techniques---input-assessment reward and gradient-free parameter learning---could in principle be applied to optimize configuration parameters in systems that serve harmful content; we recommend that practitioners apply these methods only within systems that maintain appropriate content safety guardrails.
We release all code and data under open licenses (Apache 2.0 for code, CC BY 4.0 for data) to support reproducibility and downstream research.
